\RequirePackage{plautopatch,fix-cm}
\documentclass[uplatex,dvipdfmx,a4paper]{jsarticle}
\usepackage[haranoaji]{pxchfon}
\usepackage[T1]{fontenc}
\usepackage[utf8]{inputenc}
\usepackage{lmodern}
\usepackage[dvipdfmx]{graphicx,color}
\usepackage{amsmath,amssymb,bm,braket,mathtools,mathrsfs,empheq,times,caption,cite}
\usepackage[top=25truemm,bottom=25truemm,left=20truemm,right=20truemm]{geometry} %余白の設定
\usepackage{setspace}
% 全体の行間調整
\renewcommand{\baselinestretch}{0.85}
% 図の名前を変更
\renewcommand{\figurename}{図.}
% 表の名前を変更
\renewcommand{\tablename}{Table.}
% 参考文献の名前を変更
\renewcommand{\refname}{参考文献}
% 図と図の間のスペース
\setlength\floatsep{0truemm}
% 本文と図の間のスペース
\setlength\textfloatsep{1truemm}
% \setlength\textfloatsep{0truemm}
% 本文中の図のスペース
\setlength\intextsep{0pt}
% 図とキャプションの間のスペース
\setlength\abovecaptionskip{0truemm}
% 数式上下の空白設定
\abovedisplayskip=4pt
\belowdisplayskip=4pt
% \abovedisplayskip=0pt
% \belowdisplayskip=0pt

\newcommand{\tabref}[1]{Table~\ref{#1}}
\newcommand{\equref}[1]{式~(\ref{#1})}
\newcommand{\figref}[1]{図.~\ref{#1}}

\usepackage{tcolorbox}
\usepackage{titlesec}

\titleformat*{\section}{\large\bfseries}
\newcommand{\red}[1]{\textcolor{red}{#1}}
\newcommand{\polylog}{\operatorname{polylog}}
\newcommand{\diag}{\operatorname{diag}}
% \everymath{\displaystyle} %数式を常にディスプレイスタイルにする

% sectionの前後の余白を詰める
\makeatletter
\renewcommand{\section}{%
    \@startsection{section}{1}{\z@}%
    {0.3\Cvs}{0.3\Cvs}%
    {\normalfont\normalsize\headfont\raggedright}%
}
\makeatother
% subsectionの前後の余白を詰める
\makeatletter
\renewcommand{\subsection}{%
    \@startsection{subsection}{1}{\z@}%
    {0.3\Cvs}{0.3\Cvs}%
    {\normalfont\normalsize\headfont\raggedright}%
    % {\normalfont\large\headfont\raggedright}%
}
\makeatother
% subsubsectionの前後の余白を詰める
\makeatletter
\renewcommand{\subsubsection}{%
    \@startsection{subsubsection}{1}{\z@}%
    {.25\Cvs}%
    {.25\Cvs}%
    {\headfont}
}
\makeatother

% 参考　http://www.ice.is.kit.ac.jp/~umehara/misc/comp/latex.html
\renewcommand{\topfraction}{1.2} %図表がページ上部を占める割合の上限
\renewcommand{\bottomfraction}{1.0} %図表がページ下部を占める割合の上限
\renewcommand{\dbltopfraction}{1.0} %二段組のとき段抜き図表がページ上部を占める割合の上限
\renewcommand{\textfraction}{0.01} % 文章がページを占める割合の下限
\renewcommand{\floatpagefraction}{1.0} % 図表がページを占める割合の上限
\renewcommand{\dblfloatpagefraction}{1.0} %二段組のとき段抜き図表がページを占める割合の上限
\setcounter{topnumber}{10} %ページ上部を占める図表の数の上限
\setcounter{bottomnumber}{10} %ページ下部を占める図表の数の上限
\setcounter{totalnumber}{20} %ページを占める図表の数の上限

\usepackage{url}
\usepackage[breaklinks=false,bookmarksnumbered=true,dvipdfmx]{hyperref} %ハイパーリンク
% \hypersetup{%
% bookmarksdepth=4,%
% bookmarksnumbered=true,%
% colorlinks=false,%
% urlcolor=magenta,%
% citecolor=green,%
% linkcolor=red,} %リンクの色
% 印刷時
\hypersetup{%
bookmarksdepth=4,%
bookmarksnumbered=true,%
colorlinks=true,%
urlcolor=black,%
citecolor=black,%
linkcolor=black,} %リンクの色


\begin{document}
\twocolumn[
\the\year 年度春合宿研究要旨
\hfill 指導教員 : 村松 眞由 准教授\\
\begin{spacing}{0.1}
  \centering
  \textbf{\large 宇宙輸送機非破壊検査のための赤外線応力測定と機械学習に基づくCFRP欠陥推定}
\end{spacing}
\begin{flushright}
    村松研究室 学部4年 佐藤瑞起
  \end{flushright}
]

\section{研究背景・目的}
\subsection{CFRPの宇宙輸送機器への応用}
炭素繊維強化プラスチック(carbon-fiber-reinforced plastics, CFRP)は，高比剛性・軽量性という利点から宇宙輸送機器の主要構造への適用が拡大している．特に段間構造や衛星フェアリングのような大規模・複雑形状の構造では，運用段階での健全性評価を短時間かつ高信頼に実施できる検査基盤が，再使用化と低コスト化の要求の下で求められる．一方で，CFRPは層間剝離などの内部損傷が表面観察のみでは把握しにくく，従来手法だけで効率的に健全性を保証することには限界がある\cite{nonedestruction}．

\subsection{非破壊検査と先行研究}
宇宙輸送機器の非破壊検査では打音検査や超音波測定が広く用いられてきたが，検査コストと運用時の迅速性に改善余地がある\cite{nonedestruction}．これに対して赤外線応力測定は，繰り返し荷重下の表面温度変化から表面主応力和分布(DSPSS)を取得できる非接触計測として注目されている\cite{ultra_red}．DSPSSと機械学習を組み合わせた欠陥推定の先行研究は，方法論で以下の二系統に整理できる．

第一に，畳み込みニューラルネットワーク(CNN)を用い，FEM由来DSPSSによる事前学習と赤外線応力測定データによるファインチューニングを組み合わせる二段階転移学習が提案されており，シミュレーション・実験間のドメインギャップ低減に有効である一方，対象は単純短冊形状に限定される\cite{kojima2025_transfer}．第二に，グラフニューラルネットワーク(GNN)を用いることで，CNNが前提とする格子状入力では扱いにくい不規則メッシュや穴あき形状の試験片に対しても欠陥推定が可能となるが，検証はFEMデータのみで行われており実験データとの統合は未着手である\cite{nishioka2024_bthesis,gnn}．

\subsection{研究ギャップと本研究の目的}
以上から，本研究領域には方法論的・応用的な空白が残されている．方法論的には，GNNベース欠陥推定と二段階転移学習を組み合わせる枠組みが未整備であり，両系統の長所を統合した上でのFEM-実験間ドメインギャップ低減手法と評価が確立されていない．応用的には，検証対象が単純形状・穴あき試験片までに留まり，宇宙輸送機器を意識した形状条件下での実験統合型検証が十分に行われていない．

そこで本研究では，GNNによる不規則メッシュ対応の欠陥推定モデルに，FEM由来DSPSS事前学習と赤外線応力測定データによる適応の二段階転移学習を組み合わせ，宇宙輸送機器の非破壊検査に資するCFRP欠陥推定手法の構築を目的とする．

\section{理論}
本研究では，CFRP構造のメッシュ表現と応力場情報を同時に扱うため，グラフ学習に基づく欠陥推定を想定する．
代表的なグラフ畳み込みの更新は，隣接節点の情報を集約して節点特徴を更新する構造で表される\cite{gcn}:
\begin{equation}
z_i^{(l+1)} = \phi\left(\sum_{j \in N(i) \cup \{i\}} \frac{1}{\sqrt{\deg(i)\deg(j)}}W^{(l+1)}z_j^{(l)}\right)
\end{equation}
ここで，$z_i^{(l)}$は節点特徴，$W^{(l+1)}$は学習パラメータ，$\phi$は活性化関数である．
また，有限要素解析結果の学習にグラフニューラルネットワークが有効であることが示されており，本研究のモデル設計の基盤とする\cite{gnn}．

\section{方法}
提案手法は，(i)FEM由来DSPSS，(ii)赤外線応力測定由来DSPSS，(iii)欠陥ラベルの3要素を統合して学習を行う．
まずFEMデータを用いて基礎学習を行い，次に赤外線応力測定データを用いて実験条件への適応を行う．この二段階設計により，シミュレーションデータの量的優位性と実験データの現実性を両立する．
入力は表面主応力和分布およびメッシュ関連情報，出力は欠陥位置(必要に応じて層情報を含む)とする．
評価は推定精度に加えて，ノイズ耐性，実験条件変化に対する汎化性，および実験データとの整合性を確認する．

\section{今後の予定}
短期的には，宇宙輸送機構造を対象にした解析条件の固定と，赤外線応力測定データの前処理手順の確立を行う．
次に，FEM事前学習モデルに対する実験データ適応を実施し，シミュレーション中心推定と実験統合推定の比較評価を進める．
最終的には，宇宙輸送機非破壊検査への適用を意識した，実験統合型欠陥推定手法としての妥当性を検証する．


\bibliographystyle{unsrt}
\bibliography{references}
\end{document}
